\documentclass[11pt]{article}
\usepackage[a4paper,margin=21mm]{geometry}
\usepackage[no-math]{fontspec}
\setmainfont{Latin Modern Roman}
\newfontfamily\CJKfont{Noto Serif CJK SC}
\XeTeXlinebreaklocale "zh"
\XeTeXlinebreakskip=0pt plus 1pt
\usepackage{amsmath,amssymb,amsthm,amscd,graphicx,color}
\usepackage{xurl}
\usepackage[unicode,hidelinks]{hyperref}
\allowdisplaybreaks
\setlength\emergencystretch{3em}
\numberwithin{equation}{section}

\usepackage{tikz-cd}

\numberwithin{equation}{section}

\newtheorem{Theorem}{Theorem}[section]
\newtheorem*{Theorem*}{Theorem}
\newtheorem{Corollary}[Theorem]{Corollary}
\newtheorem{Lemma}[Theorem]{Lemma}
\newtheorem{Proposition}[Theorem]{Proposition}
 { \theoremstyle{definition}
\newtheorem{Definition}[Theorem]{Definition}
\newtheorem{Note}[Theorem]{Note}
\newtheorem{Example}[Theorem]{Example}
\newtheorem{Remark}[Theorem]{Remark} }

\DeclareMathOperator{\Aut}{Aut}
\DeclareMathOperator{\saff}{SAff}
\DeclareMathOperator{\Span}{Span}
\DeclareMathOperator{\der}{Der}
\DeclareMathOperator{\stab}{Stab}
\DeclareMathOperator{\fix}{Fix}
\DeclareMathOperator{\sgn}{sign}
\DeclareMathOperator{\Imag}{Im}
\DeclareMathOperator{\Real}{Re}
\DeclareMathOperator{\End}{End}
\DeclareMathOperator{\rk}{\text{rk}}
\DeclareMathOperator{\Hom}{Hom}
\DeclareMathOperator{\Ad}{Ad}
\DeclareMathOperator{\ad}{ad}
\DeclareMathOperator{\pr}{Pr}
\DeclareMathOperator{\Coad}{Ad^*}
\DeclareMathOperator{\Gr}{Gr}
\DeclareMathOperator{\spoon}{\text{Span}}
\DeclareMathOperator{\Tr}{Trace}
\DeclareMathOperator{\aff}{Aff}
\DeclareMathOperator{\affo}{\mathfrak{aff}}
\DeclareMathOperator{\coad}{ad*}
\DeclareMathOperator{\crit}{Crit}
\DeclareMathOperator{\F}{\mathcal{F}}
\newcommand{\orb}{\mathcal{O}}
\newcommand{\Left}{\mathcal{L}}
\newcommand{\Right}{\mathcal{R}}
\newcommand{\R}{\mathbb{R}}
\newcommand{\C}{\mathbb{C}}
\newcommand{\CP}{\mathbb{C}P}
\newcommand{\D}{\mathcal{D}}
\newcommand{\CS}{\mathbb{CS}}
\newcommand{\RP}{\mathbb{R}P}
\newcommand{\Q}{\mathbb{H}}
\newcommand{\Symp}{\text{Symp}}
\newcommand{\symp}{\text{symp}}
\newcommand{\Id}{{\rm Id}}

\makeatletter\newlabel{list_of_RST}{{3.6}{}{Original source locator}{}{}}\makeatother
\begin{document}
\begin{center}\Large Holomorphic symplectic Grassmannian and coadjoint-orbit models of the hyperkähler U(m) quotient at (i Id\_m,0,0) for m<=n\end{center}
\noindent\textbf{Stable ID:} 1054\quad\textbf{Author:} Philip Arathoon; Marine Fontaine\par
\noindent\textbf{Work:} Real Forms of Holomorphic Hamiltonian Systems\par
\noindent\textbf{Source:} \url{https://sigma-journal.com/2024/114/}\par
\noindent\textbf{Version:} Official SIGMA 20 (2024) article 114, DOI 10.3842/SIGMA.2024.114; journal PDF and native TeX acquired 2026-09-30, exact bytes pinned by SHA256\par
\noindent\textbf{Rights:} CC BY 4.0. Authors retain copyright. \url{https://creativecommons.org/licenses/by/4.0/}. Reuse and typesetting adaptation; original language and mathematical content preserved.\par
\noindent\textbf{Difficulty (prior selection preserved):} {\CJKfont H2: The author constructs the global quotient identification by proving that every complex matrix orbit in the specified open set meets the real moment level in exactly one unitary orbit. Singular-value reduction provides the existence and uniqueness argument, while three explicit quaternionic complex-coordinate descriptions identify the rotated models and their symplectic forms. The global real-versus-complex quotient step is the substantive obligation beyond a local coordinate computation.}\par
\medskip\noindent\textbf{Editorial boundary note (not author text):} Complete original Theorem4.4, all three holomorphic symplectic model identifications, and full existence/uniqueness of the complex-orbit intersection with a single unitary orbit. The full trace pairing, canonical symplectic form, group actions and momentum maps, coadjoint-orbit definition and KKS form, Hermitian and quaternionic metrics, all three quaternionic coordinate decompositions and their author proof, regular level and rank hypotheses are included. The supplied coordinate/metric verification and rotated-complex-structure repetition retain the author wording and explicit inputs, without editor verification or new proof. The prior dual-pair fact remains the exact author reference. This substantial global model-classification result is one selected problem unit; the coordinate lemma is supporting context, and the later diffeomorphism/brane/mechanical outcomes are unselected. Literal source notation and language are preserved.\par
\noindent\textbf{External original reference list\_of\_RST:} Original matrix real-structure involutions, source367–376, original PDF10. Referenced only by the retained metric remark regarding isometry compatibility; not an omitted input or proof of the selected quotient-model theorem.\par
\medskip\hrule\medskip\noindent\textbf{Author text begins below.}\par
\setcounter{section}{3}\setcounter{subsection}{1}\setcounter{equation}{3}\setcounter{Theorem}{3}
% Original source lines 323--354
\subsection{A coadjoint orbit example}\label{theexample}
Consider the space $\Hom(\C^m,\C^n)$ of complex $n\times m$ matrices $Q$. The dual space is identified with matrices $P$ in $\Hom(\C^n,\C^m)$ via the trace pairing $\Tr(PQ)$. The holomorphic cotangent bundle $T^*_{1,0}\Hom(\C^m,\C^n)$ is then identified with the set of such pairs $(Q,P)$. The holomorphic symplectic form is
\begin{equation*}%\label{symp_form_on_v_space}
	\Omega((Q_1,P_1),(Q_2,P_2))=\Tr(P_2Q_1-P_1Q_2).
\end{equation*}
The groups ${\rm GL}_m\C$ and ${\rm GL}_n\C$ act symplectically by $\bigl(Qg^{-1},gP\bigr)$ and $\bigl(hQ,Ph^{-1}\bigr)$ with momentum maps
\begin{equation*}	%\label{momentum_maps}
	\mu_m(Q,P)=PQ,\qquad\text{and}\qquad\mu_n(Q,P)=QP.
\end{equation*}
Here we have identified $\mathfrak{gl}_n$ and $\mathfrak{gl}_m$ with their respective duals using the trace form.

Now suppose $m\le n$ and let $M$ be the open subset of the cotangent bundle where both $Q$ and $P$ have maximal rank. One can show (see \cite{dualpairs}) that $\mu_n$ is an orbit-map for the ${\rm GL}_m\C$-action on $M$ and that the fibres of $\mu_m$ are orbits of ${\rm GL}_n\C$. As the momentum map is Poisson we can identify the Poisson reduced space $M/{\rm GL}_m\C$ with $\mu_n(M)\subset\mathfrak{gl}_n\C^*$. For $\zeta\in\C^\times$, the orbit-reduced space $\mu_m^{-1}(\zeta\cdot\Id_m)/{\rm GL}_m\C$ is identified with the coadjoint orbit
\begin{equation}
	\label{coad_orbit}
	{\rm Orb}(\zeta)=\{\xi=QP\mid (Q,P)\in M,\,PQ=\zeta\cdot\Id_m\}\subset\mathfrak{gl}_n\C^*.
\end{equation}

Geometrically, $\xi$ determines a decomposition $\C^n=\Sigma^{(m)}\oplus \Lambda^{(n-m)}$, where ${\Sigma=\Imag Q}$ and $\Lambda=\ker P$. Conversely, such a decomposition determines an element $\xi$ where $\ker\xi=\Lambda$ and ${\xi|_\Sigma=\zeta\cdot\Id_\Sigma}$. Naturally, we see that the orbit is diffeomorphic to the symmetric space $\D_\C$ given~by
\begin{Definition}
	For $\mathbb{F}=\R,\C,\Q$, the \emph{space of decompositions} $\D_\mathbb{F}$ is the set
	\begin{equation}
		\label{elegant_A+B}
		\D_\mathbb{F}=(\Gr_\mathbb{F}(m;n)\times \Gr_\mathbb{F}(n-m;n))\setminus\Delta,
	\end{equation}
	where $\Delta$ is the closed subset of pairs $(\Sigma,\Lambda)$ with $\Sigma\cap\Lambda\ne\{0\}$.
\end{Definition}
\begin{Remark}\label{elegant_rmk}
	The symplectic form on $\D_{\C}$ admits a rather elegant description. Recall that tangent vectors to an element $\Sigma$ of the Grassmannian may be identified with linear maps ${\Sigma\rightarrow\C^n/\Sigma}$. Since $(\Sigma,\Lambda)$ determines a decomposition of $\C^n$, the tangent space to $\D_{\C}$ at this point may be identified with pairs of maps $(\alpha\colon\Sigma\rightarrow\Lambda,\,\beta\colon \Lambda\rightarrow\Sigma)$. The KKS form on ${\rm Orb}(\zeta)$ can be shown~to~be
	\[
		\Omega_{\rm KKS}((\alpha_1,\beta_1),(\alpha_2,\beta_2))=\zeta\Tr(\beta_2\alpha_1-\beta_1\alpha_2).
	\]
\end{Remark}

\setcounter{subsection}{3}\setcounter{Theorem}{7}\setcounter{equation}{6}
% Original source lines 394--405
\begin{Remark}%\label{kahler metric}
	Let $\langle x,y\rangle=x^\dagger y$ denote the ordinary Hermitian forms on $\C^n$ and $\C^m$. This allows us to equip $T^*_{1,0}\Hom(\C^m,\C^n)$ with a K\"{a}hler metric
	\begin{equation}\label{metric}
		g((Q_1,P_1),(Q_2,P_2))=\frac{1}{2}\Tr\bigl(Q_1Q_2^\dagger+Q_2Q_1^\dagger\bigr)+\frac{1}{2}\Tr\bigl(P_1P_2^\dagger+P_2P_1^\dagger\bigr).
	\end{equation}
	Looking ahead, we would like the involutions in \eqref{list_of_RST} to be isometries of this metric. To ensure this, we will from now impose the following:
	\begin{itemize}\itemsep=0pt
		\item Compatibility between the real structure on $\C^n$ with the standard Hermitian form $\langle \,,\,\rangle$, the~Hermitian form $(\,,\,)$, and the quaternionic structures $\mathbb{J}$, in the sense that: $\langle\overline{x},\overline{y}\rangle=\overline{\langle x, y\rangle}$, $(\overline{x},\overline{y})=\overline{(x,y)}$, and $\mathbb{J}(\overline{x})=\overline{\mathbb{J}(x)}$.
		\item Compatibility between the Hermitian form $(\,,\,)$ defined by $\Theta$ and quaternionic structures~$\mathbb{J}$ with $\langle\,,\,\rangle$ in the sense that: $\langle \Theta x,\Theta y\rangle=\langle x,y\rangle$ and $\langle\mathbb{J}(x),\mathbb{J}(y)\rangle=\langle x,y\rangle$.
	\end{itemize}
	In plainer terms, if $x\mapsto\overline{x}$ is ordinary complex conjugation on $\C^n$, we require $\Theta$ to be a real symmetric matrix with $\Theta^2=\Id_n$, and $\mathbb{J}(x)=\mathcal{J}\overline{x}$ for $\mathcal{J}$ a real skew-symmetric matrix.
\end{Remark}

\setcounter{section}{3}
% Original source lines 407--414
\section{Branes}
\subsection{Hyperk\"{a}hler geometry}
A hyperk\"{a}hler manifold is a Riemannian manifold $(M,g)$ equipped with three complex structures, $I$, $J$, and $K$, which satisfy the usual quaternionic relations, and for which $(M,g)$ is K\"{a}hler with respect to each of them. A hyperk\"{a}hler manifold defines a holomorphic symplectic form on $M$ with respect to each complex structure. If we denote the K\"{a}hler forms by
\[
\omega_1(X,Y)=g(I(X),Y),\qquad\omega_2(X,Y)=g(J(X),Y),\qquad\omega_3(X,Y)=g(K(X),Y),
\]
then $\Omega_1=\omega_2+{\rm i}\omega_3$, $\Omega_2=\omega_3+{\rm i}\omega_1$, and $\Omega_3=\omega_1+{\rm i}\omega_2$ each define holomorphic symplectic forms on $M$ with respect to the complex structures, $I$, $J$, and $K$, respectively.


\setcounter{Theorem}{2}
% Original source lines 428--483
\subsection{A hyperk\"{a}hler structure on the coadjoint orbit}
Consider the quaternionic spaces $\Q^m$ and $\Q^n$ viewed as right $\Q$-modules. Then the space $\Hom(\Q^m,\Q^n)$ of $n\times m$ quaternionic matrices is itself a left $\Q$-module, where we denote left multiplication by the quaternions ${\rm i}$, ${\rm j}$, ${\rm k}$
with the operators $I$, $J$, $K$. We can further equip this space with a metric%
\begin{equation}\label{quat_metric}
	g(A,B)=\frac{1}{2}\Tr\bigl(AB^\dagger+BA^\dagger\bigr),
\end{equation}
where $A^\dagger$ is the quaternionic conjugate-transpose. The tuple $(g,I,J,K)$ defines a hyperk\"{a}hler structure on $\Hom(\Q^m,\Q^n)$.
Observe that right multiplication by the compact symplectic group ${\rm Sp}(m)\subset {\rm GL}_m\Q$ preserves the hyperk\"{a}hler structure. The following lemma will facilitate us in switching between different complex structures on this space.
\begin{Lemma}\label{facilitate}
	By writing elements $A$ of $\Hom(\Q^m,\Q^n)$ as
	\begin{equation*}
		A=\begin{cases}
			\displaystyle Q+{\rm j}P^\dagger,\\[1mm]
			\displaystyle\frac{Y^\dagger-X}{\sqrt{2}}-{\rm j}\biggl(\frac{{\rm i}\overline{X}+{\rm i}Y^{\mathsf{T}}}{\sqrt{2}}\biggr),\vspace{1mm}\\
			\displaystyle\frac{U+{\rm i}V^\dagger}{\sqrt{2}}+{\rm j}\biggl(\frac{{\rm i}V^\dagger-U}{\sqrt{2}}\biggr)
		\end{cases}
	\end{equation*}
	for $(Q,P)$, $(X,Y)$, and $(U,V)$ pairs of matrices in $\Hom(\C^m,\C^n)\times\Hom(\C^n,\C^m)$, we establish isometric holomorphic symplectomorphisms between $T^*_{1,0}\Hom(\C^m,\C^n)$ and $\Hom(\Q^m,\Q^n)$ for the structures $(I,\Omega_1)$, $(J,\Omega_2)$, and $(K,\Omega_3)$, respectively. Furthermore, these intertwine the ${\rm U}(m)\subset {\rm Sp}(m)$-action with the ${\rm U}(m)\subset {\rm GL}_m\C$-action on $T^*_{1,0}\Hom(\C^m,\C^n)$.
\end{Lemma}
\begin{proof}
	It is straightforward to verify that the metric in \eqref{metric} is pulled back to \eqref{quat_metric}. One can then show that for the three choices of variables above, the roles of $I$, $J$, and $K$ are cyclically permuted, and thus, so are the holomorphic symplectic forms $\Omega_1$, $\Omega_2$, and $\Omega_3$.
\end{proof}

From now on, suppose $m\le n$ and let $M$ denote the open subset of $\Hom(\Q^m,\Q^n)$ consisting of all elements $A=Q+{\rm j}P^\dagger$ for which $Q$ has maximal rank. The ${\rm U}(m)$-action on $M$ admits a~triple of momentum maps $\mu_k\colon M\rightarrow\mathfrak{u}(m)^*$ for each K\"{a}hler form $\omega_k$. This allows us to introduce the hyperk\"{a}hler quotient
\[
\tilde{M}=\bigl(\mu_1^{-1}({\rm i}\cdot\Id_m)\cap\mu_2^{-1}(0)\cap\mu_3^{-1}(0)\bigr)/{\rm U}(m).
\]
\begin{Theorem}\label{HK_red_thm}
	Let $\bigl(\tilde{M},g,I,J,K\bigr)$ be the hyperk\"{a}hler reduced space for the action of ${\rm U}(m)\subset {\rm Sp}(m)$ on $M\subset\Hom(\Q^m,\Q^n)$ at the regular value $({\rm i}\cdot\Id_m,0,0)\in\mathfrak{u}(m)^*\otimes\R^3$. We have the~following holomorphic symplectomorphisms:
	\begin{itemize}\itemsep=0pt
		\item $\bigl(\tilde{M}, I, \Omega_1\bigr)\cong (T^*_{1,0}\Gr_\C, {\rm i}, \Omega_{\rm can})$,
		\item $\bigl(\tilde{M},J,\Omega_2\bigr)\cong({\rm Orb}(-1), {\rm i}, \Omega_{\rm KKS})$,
		\item $\bigl(\tilde{M},K,\Omega_3\bigr)\cong({\rm Orb}({\rm i}), {\rm i}, \Omega_{\rm KKS})$.
	\end{itemize}
	Here ${\rm Orb}(\zeta)$ is the coadjoint orbit in $\mathfrak{gl}_n\C^*$ equipped with the Kostant--Kirilov--Souriau form $\Omega_{\rm KKS}$ and $\Gr_\C$ is the Grassmannian of complex $m$-dimensional subspaces in $\C^n$.
\end{Theorem}

\begin{proof}
	For the distinguished complex structure $I$, we use Lemma~\ref{facilitate} to identify $\Hom(\Q^m,\Q^n)$ with ${T^*_{1,0}\Hom(\C^m,\C^n)}$. The ${\rm U}(m)$-action admits a holomorphic extension to the ${\rm GL}_m\C$-action demonstrated in Section~\ref{theexample} with momentum map $\mu_m=\mu_2+{\rm i}\mu_3$ given by $\mu_m(Q,P)=PQ$.
	
	In this basis $\mu_1(Q,P)={\rm i}\bigl(Q^\dagger Q-PP^\dagger\bigr)$. We claim that each ${\rm GL}_m\C$-orbit in $M$ intersects $\mu_1^{-1}({\rm i}\cdot\Id_m)$ precisely in a single orbit of ${\rm U}(m)$. To show this, consider how the action of $g$ in ${\rm GL}_m\C$ sends $Q^\dagger Q-PP^\dagger$ to
	\begin{equation}\label{svd}
		g^{-\dagger}Q^\dagger Qg^{-1}-gPP^\dagger g^{\dagger}.
	\end{equation}
	If $Q$ has maximal rank, we may act by the appropriate $g$ to assume that $Q^\dagger Q$ is the identity. If~we~then use the singular value decomposition $g=uDv^\dagger$ and suppose that $v^\dagger PP^\dagger v$ is a~diagonal matrix $\Lambda$ with non-negative diagonal entries, then \eqref{svd} becomes $u\bigl(D^{-\dagger}D^{-1}-D\Lambda D^\dagger\bigr)u^\dagger$. We~can always find a $D$ for which this equals the identity, and hence, have established that the intersection is non-empty.
	
	To show that this intersection is a single ${\rm U}(m)$-orbit, let $\mu_1(Q,P)={\rm i}\cdot\Id_m$ and suppose that~\eqref{svd} is equal to $\Id_m$. By again writing $g$ as $uDv^\dagger$, we find from substituting ${Q^\dagger Q\!=\Id_m\!+\!PP^\dagger}$ into \eqref{svd} that we must have $D^{-\dagger}(\Id_m+S)D^{-1}=\Id_m+DSD^\dagger$, where $S=v^\dagger PP^\dagger v$. This can only hold when $D$ is the identity, and hence, $g$ is unitary.
	
	From this claim, it follows that $\bigl(\tilde{M},I,\Omega_1\bigr)$ may be identified with the holomorphic symplectic reduced space
	\begin{equation*}%\label{QP_quotient}
		\mu_m^{-1}(0)/{\rm GL}_m\C.
	\end{equation*}
	There is a well-defined identification between this quotient and $T^*_{1,0}\Gr_\C$ which sends the equivalence class $[(Q,P)]$ to the plane $\Pi=\Imag Q$ and covector $QP$, interpreted as a map $\C^n/\Pi\rightarrow\C^n$. Observe that the canonical one-form on the cotangent bundle pulls back under this identification to the canonical one-form on $T^*_{1,0}\Hom(\C^m,\C^n)$.
	
	For a different choice of distinguished complex structure, we apply the same argument. For~$J$, we must consider $\mu_m^{-1}(-\Id_m)$ for $\mu_m=\mu_3+{\rm i}\mu_1$, and for $K$ we have $\mu_m^{-1}({\rm i}\cdot\Id_m)$ where ${\mu_m=\mu_1+{\rm i}\mu_2}$. In both cases, the momentum conditions imply that the matrices $(X,Y)$ and~$(U,V)$ all have maximal rank, and a slight alteration to the claim above establishes that the quotients may be identified with the holomorphic symplectic reduced spaces in \eqref{coad_orbit}.
\end{proof}
\begin{thebibliography}{99}
\bibitem[29]{dualpairs}
Skerritt P., Vizman C., Dual pairs for matrix groups, \href{https://doi.org/10.3934/jgm.2019014}{\textit{J.~Geom. Mech.}}
 \textbf{11} (2019), 255--275, \href{https://arxiv.org/abs/1805.01519}{arXiv:1805.01519}.

\end{thebibliography}
\end{document}
